% GENERATED FILE. Edit canonical lessons, then run npm run generate:books.
% canonical-source-hash: fnv1a64:a7de2602733df0b6
% canonical-lessons: RU-C525-poet, RU-C525-stikhotvoreniye, RU-C525-rifma, RU-C525-rifmovat, RU-C525-naizust

\chapter{Poems and Rhymes}
\label{ch:ru-poems-and-rhymes}
\hlchaptermodality{525}

\begin{chapteropening}
\textbf{By the end of this chapter:} \emph{I can talk about poems and learning them by heart.}

Complete the last lesson of chapter 525: I can talk about poems and learning them by heart.
\end{chapteropening}

\section[\texorpdfstring{\ru{поэт} (poét)}{поэт (poét)}]{\ru{поэт} (poét) — a poet}

\label{lesson:RU-C525-poet}



\begin{quote}
Before the new one: say the Russian for a novella, then the Russian for memoirs.
\end{quote}

\subsection*{You'll want to know: \ru{поэт}}
\textbf{\ru{поэт}} (\emph{poét}) — \textquotedblleft{}a poet\textquotedblright{}.

\begin{quote}
\ru{Этот} \ru{поэт} \ru{писал} \ru{о} \ru{любви} \ru{и} \ru{природе}.
\end{quote}

\begin{quote}
\emph{Étot poét pisál o lyubví i priróde.}
\end{quote}

\begin{quote}
This poet wrote about love and nature.
\end{quote}

Masculine; feminine \ru{поэтесса}, though many women prefer \emph{\ru{поэт}}.

\begin{cousinweb}
\emph{\ru{поэт}} was borrowed through German \emph{Poet} from Latin \emph{poēta}, which comes from Greek \emph{poiēt\'{\={e}}s} 'a maker, a poet' (EtymDB).
\end{cousinweb}

\begin{grammarlens}[title={What you've built}]
One more word for telling, explaining and describing at B1.
\end{grammarlens}

\subsection*{Your turn}
\begin{itemize}
\raggedright
  \item \emph{Say it:} \emph{poét}
  \item \emph{Say it:} \emph{poét}, once more
  \item \emph{Say it:} a sentence of your own with \emph{poét}
  \item \emph{Recall:} say the Russian for a novella, then the Russian for memoirs, then say \emph{poét} again
\end{itemize}

\begin{tcolorbox}[breakable,colback=teal!4,colframe=teal!35!black,title={Before you move on}]
What does \ru{поэт} mean? (\textquotedblleft{}A poet\textquotedblright{}.) Say it once more.
\end{tcolorbox}
\section[\texorpdfstring{\ru{стихотворение} (stikhotvoréniye)}{стихотворение (stikhotvoréniye)}]{\ru{стихотворение} (stikhotvoréniye) — a poem}

\label{lesson:RU-C525-stikhotvoreniye}



\begin{quote}
Before the new one: say the Russian for memoirs, then the Russian for a poet.
\end{quote}

\subsection*{You'll want to know: \ru{стихотворение}}
\textbf{\ru{стихотворение}} (\emph{stikhotvoréniye}) — \textquotedblleft{}a poem\textquotedblright{}.

\begin{quote}
\ru{В} \ru{школе} \ru{мы} \ru{учили} \ru{это} \ru{стихотворение}.
\end{quote}

\begin{quote}
\emph{V shkóle my uchíli éto stikhotvoréniye.}
\end{quote}

\begin{quote}
We learnt this poem at school.
\end{quote}

Neuter. \emph{\ru{Стихи}} — verse, poems.

\begin{cousinweb}
\emph{\ru{стихотворение}} is a compound of \emph{\ru{стих}} 'a line of verse' and \emph{\ru{творить}} 'to create'; EtymDB has no entry.
\end{cousinweb}

\begin{grammarlens}[title={What you've built}]
One more word for telling, explaining and describing at B1.
\end{grammarlens}

\subsection*{Your turn}
\begin{itemize}
\raggedright
  \item \emph{Say it:} \emph{stikhotvoréniye}
  \item \emph{Say it:} \emph{stikhotvoréniye}, once more
  \item \emph{Say it:} a sentence of your own with \emph{stikhotvoréniye}
  \item \emph{Recall:} say the Russian for memoirs, then the Russian for a poet, then say \emph{stikhotvoréniye} again
\end{itemize}

\begin{tcolorbox}[breakable,colback=teal!4,colframe=teal!35!black,title={Before you move on}]
What does \ru{стихотворение} mean? (\textquotedblleft{}A poem\textquotedblright{}.) Say it once more.
\end{tcolorbox}
\section[\texorpdfstring{\ru{рифма} (rífma)}{рифма (rífma)}]{\ru{рифма} (rífma) — a rhyme}

\label{lesson:RU-C525-rifma}



\begin{quote}
Before the new one: say the Russian for a poet, then the Russian for a poem.
\end{quote}

\subsection*{You'll want to know: \ru{рифма}}
\textbf{\ru{рифма}} (\emph{rífma}) — \textquotedblleft{}a rhyme\textquotedblright{}.

\begin{quote}
\ru{В} \ru{современных} \ru{стихах} \ru{часто} \ru{нет} \ru{рифмы}.
\end{quote}

\begin{quote}
\emph{V sovreménnykh stikhákh chásto net rífmy.}
\end{quote}

\begin{quote}
Modern poems often have no rhyme.
\end{quote}

Feminine.

\begin{cousinweb}
\emph{\ru{рифма}} is usually traced through German or Polish to French \emph{rime}; EtymDB has no entry.
\end{cousinweb}

\begin{grammarlens}[title={What you've built}]
One more word for telling, explaining and describing at B1.
\end{grammarlens}

\subsection*{Your turn}
\begin{itemize}
\raggedright
  \item \emph{Say it:} \emph{rífma}
  \item \emph{Say it:} \emph{rífma}, once more
  \item \emph{Say it:} a sentence of your own with \emph{rífma}
  \item \emph{Recall:} say the Russian for a poet, then the Russian for a poem, then say \emph{rífma} again
\end{itemize}

\begin{tcolorbox}[breakable,colback=teal!4,colframe=teal!35!black,title={Before you move on}]
What does \ru{рифма} mean? (\textquotedblleft{}A rhyme\textquotedblright{}.) Say it once more.
\end{tcolorbox}
\section[\texorpdfstring{\ru{рифмовать} (rifmovát)}{рифмовать (rifmovát)}]{\ru{рифмовать} (rifmovát) — to rhyme}

\label{lesson:RU-C525-rifmovat}



\begin{quote}
Before the new one: say the Russian for a poem, then the Russian for a rhyme.
\end{quote}

\subsection*{You'll want to know: \ru{рифмовать}}
\textbf{\ru{рифмовать}} (\emph{rifmovát}) — \textquotedblleft{}to rhyme\textquotedblright{}.

\begin{quote}
\ru{Мой} \ru{сын} \ru{любит} \ru{рифмовать} \ru{слова}.
\end{quote}

\begin{quote}
\emph{Moy syn lyúbit rifmovát slová.}
\end{quote}

\begin{quote}
My son likes making words rhyme.
\end{quote}

Imperfective; present \ru{рифмую}. \emph{\ru{Рифмоваться}} — to rhyme (of words).

\begin{cousinweb}
\emph{\ru{рифмовать}} is built on \emph{\ru{рифма}} 'a rhyme' with \emph{-\ru{овать}}; EtymDB has no entry, and \emph{\ru{рифма}} is usually traced to French \emph{rime}.
\end{cousinweb}

\begin{grammarlens}[title={What you've built}]
One more word for telling, explaining and describing at B1.
\end{grammarlens}

\subsection*{Your turn}
\begin{itemize}
\raggedright
  \item \emph{Say it:} \emph{rifmovát}
  \item \emph{Say it:} \emph{rifmovát}, once more
  \item \emph{Say it:} a sentence of your own with \emph{rifmovát}
  \item \emph{Recall:} say the Russian for a poem, then the Russian for a rhyme, then say \emph{rifmovát} again
\end{itemize}

\begin{tcolorbox}[breakable,colback=teal!4,colframe=teal!35!black,title={Before you move on}]
What does \ru{рифмовать} mean? (\textquotedblleft{}To rhyme\textquotedblright{}.) Say it once more.
\end{tcolorbox}
\section[\texorpdfstring{\ru{наизусть} (naizúst)}{наизусть (naizúst)}]{\ru{наизусть} (naizúst) — by heart}

\label{lesson:RU-C525-naizust}



\begin{quote}
Before the new one: say the Russian for a rhyme, then the Russian for to rhyme.
\end{quote}

\subsection*{You'll want to know: \ru{наизусть}}
\textbf{\ru{наизусть}} (\emph{naizúst}) — \textquotedblleft{}by heart\textquotedblright{}.

\begin{quote}
\ru{Бабушка} \ru{знает} \ru{наизусть} \ru{много} \ru{стихов}.
\end{quote}

\begin{quote}
\emph{Bábushka znáyet naizúst mnógo stikhóv.}
\end{quote}

\begin{quote}
Grandma knows a lot of poems by heart.
\end{quote}

Adverb. \emph{\ru{Учить} \ru{наизусть}} — to learn by heart.

\begin{cousinweb}
\emph{\ru{наизусть}} is usually explained as \emph{\ru{на}}, \emph{\ru{из}} and the old word \emph{\ru{уста}} 'lips, mouth', so it means 'from the mouth'; EtymDB has no entry for \emph{\ru{наизусть}}, and \emph{\ru{уста}} is inherited from Proto-Slavic \emph{usta} (EtymDB).
\end{cousinweb}

\begin{grammarlens}[title={What you've built}]
One more word for telling, explaining and describing at B1.
\end{grammarlens}

\subsection*{Your turn}
\begin{itemize}
\raggedright
  \item \emph{Say it:} \emph{naizúst}
  \item \emph{Say it:} \emph{naizúst}, once more
  \item \emph{Say it:} a sentence of your own with \emph{naizúst}
  \item \emph{Recall:} say the Russian for a rhyme, then the Russian for to rhyme, then say \emph{naizúst} again
\end{itemize}

\begin{tcolorbox}[breakable,colback=teal!4,colframe=teal!35!black,title={Before you move on}]
What does \ru{наизусть} mean? (\textquotedblleft{}By heart\textquotedblright{}.) Say it once more.
\end{tcolorbox}
